\documentclass[12pt]{article}
\usepackage[utf8]{inputenc}
\usepackage[T1]{fontenc}
\usepackage{hyperref}
\title{Tandia Language Resources for K-12 Curriculum Development}
\author{Compiled by Claude Sonnet 5 (claude-sonnet-5) for Andr\'as Kornai}
\date{}
\begin{document}
\maketitle

\section{Introduction}
Tandia (ISO 639-3 tni; Glottocode tand1253) is an Austronesian language of the Wandamen Peninsula in West Papua, Indonesia. It belongs to the South Halmahera--West New Guinea branch. In 1991 only two speakers were reported. Wikipedia states that a survey found it extinct by 2024, while ASJP lists it as extinct since 2000. Speakers shifted to Wandamen (Wamesa). Very little is documented, and the searches found no teaching materials.

\section{Primary Educational Resources}
None found. A heritage curriculum would have to be built from the scant documentation that exists.

\section{Linguistic and Reference Materials}
\begin{itemize}
\item Wikipedia, Tandia language: \url{https://en.wikipedia.org/wiki/Tandia_language}
\item Glottolog entry: \url{https://glottolog.org/resource/languoid/id/tand1253.bigmap.html}
\item ASJP word list: \url{https://asjp.clld.org/languages/TANDIA}
\item TransNewGuinea.org language page: \url{https://transnewguinea.org/language/tandia}
\end{itemize}

\section{Religious and Community Resources}
None found.

\section{Digital and Audio Resources}
No audio or video was found. A Verbix page on endangered languages lists Tandia, but I did not verify its content: \url{https://docs.verbix.com/EndangeredLanguages/Tandia}

\section{Curriculum Development Context}
The language is dormant or extinct, with no first-language speakers. A revitalization curriculum would rest on archival word lists. It could be developed alongside the neighbouring Wandamen (Wamesa) community, to which the Tandia speakers shifted. No named local contacts were found. The honest summary is that little exists.

\end{document}